% Part: incompleteness
% Chapter: representability-in-q
% Section: composition-representable

\documentclass[../../../include/open-logic-section]{subfiles}

\begin{document}

\olfileid{inc}{req}{cmp}
\olsection{સંયોજન~$\Th{Q}$માં નિરૂપ્ય છે}

માનો કે $h$ આ પ્રમાણે
વ્યાખ્યાયિત છે:
\[
h(x_0,\dots,x_{l-1}) = f(g_0(x_0,\dots,x_{l-1}), \dots,
g_{k-1}(x_0,\dots,x_{l-1})).
\]
અહીં $f$ અને અનુક્રમે $g_0$,
\dots,~$g_{k-1}$ને નિરૂપતાં
!!{formula}s $!A_f, !A_{g_0}, \dots,
!A_{g_{k-1}}$ આપણે અગાઉથી શોધ્યાં
છે. હવે~$h$ને નિરૂપતું
!!a{formula}~$!A_h$ શોધવાનું છે.

પહેલાં બધાં વિધેયો $1$-સ્થાની
હોય એવો સરળ કિસ્સો લો, એટલે
$h(x) = f(g(x))$ વિચારો.
$!A_f(y, z)$ એ~$f$ને અને
$!A_g(x, y)$ એ~$g$ને નિરૂપે,
તો~$h$ને નિરૂપતું
!!a{formula}~$!A_h(x, z)$ જોઈએ.
ધ્યાન આપો કે $h(x) = z$
ત્યારે અને માત્ર ત્યારે જ એવું~$y$
છે કે $z = f(y)$ અને
$y = g(x)$ બંને સાચાં છે.
($h(x) = z$ હોય, તો $g(x)$
આવું~$y$ છે; અને આવું $y$
હોય, તો $y = g(x)$ તથા
$z = f(y)$ હોવાથી $z
= f(g(x))$.) તેથી
$\lexists[y][(!A_g(x, y) \land !A_f(y,
  z))]$ એ~$!A_h(x, z)$ માટે
યોગ્ય ઉમેદવાર લાગે છે.
હવે સંબંધિત !!{formula}s
$\Th{Q}$માં સાબિત થાય છે
તેની ખાતરી કરવી બાકી છે.

\begin{prop}
\ollabel{prop:rep1}
$h(n) = m$ હોય, તો
$\Th{Q} \Proves !A_h(\num{n}, \num{m})$.
\end{prop}

\begin{proof}
માનો કે $h(n) = m$, એટલે કે
$f(g(n)) = m$. $k = g(n)$ લો.
તો
\begin{align*}
  \Th{Q} & \Proves !A_g(\num{n}, \num{k})
  \intertext{કારણ કે $!A_g$ એ~$g$ને નિરૂપે છે, અને}
  \Th{Q} & \Proves !A_f(\num{k}, \num{m})
  \intertext{કારણ કે $!A_f$ એ~$f$ને નિરૂપે છે. તેથી}
  \Th{Q} & \Proves !A_g(\num{n}, \num{k}) \land !A_f(\num{k}, \num{m})
  \intertext{અને પરિણામે}
  \Th{Q} & \Proves \lexists[y][(!A_g(\num{n}, y) \land !A_f(y, \num{m}))],
\end{align*}
એટલે કે $\Th{Q} \Proves
!A_h(\num{n}, \num{m})$.
\end{proof}

\begin{prop}
\ollabel{prop:rep2}
$h(n) = m$ હોય, તો
$\Th{Q} \Proves \lforall[z][(!A_h(\num{n}, z) \lif
  z = \num{m})]$.
\end{prop}

\begin{proof}
માનો કે $h(n) = m$, એટલે કે
$f(g(n)) = m$. $k = g(n)$ લો.
તો
\begin{align*}
  \Th{Q} & \Proves \lforall[y][(!A_g(\num{n}, y) \lif \eq[y][\num{k}])]
  \intertext{કારણ કે $!A_g$ એ~$g$ને નિરૂપે છે, અને}
  \Th{Q} & \Proves \lforall[z][(!A_f(\num{k}, z) \lif \eq[z][\num{m}])]
  \intertext{કારણ કે $!A_f$ એ~$f$ને નિરૂપે છે. થોડો તર્ક વાપરીને આ પણ બતાવી શકીએ:}
  \Th{Q} & \Proves \lforall[z][(\lexists[y][(!A_g(\num{n}, y) \land
      !A_f(y, z))] \lif \eq[z][\num{m}])].
\end{align*}
એટલે કે $\Th{Q} \Proves
\lforall[y][(!A_h(\num n, y) \lif \eq[y][\num m])]$.
\end{proof}

$f$ અને~$g_i$ની સ્થાનસંખ્યા
$1$ કરતાં વધુ હોય એવા જટિલ
કિસ્સામાં પણ આ જ વિચાર કામ કરે છે.

\begin{prop}
\ollabel{prop:rep-composition}
$!A_f(y_0, \dots, y_{k-1}, z)$
એ~$f(y_0, \dots, y_{k-1})$ને
$\Th{Q}$માં નિરૂપે, અને
$!A_{g_i}(x_0, \dots, x_{l-1}, y)$
એ $g_i(x_0, \dots, x_{l-1})$ને
$\Th{Q}$માં નિરૂપે, તો આ સૂત્ર
\begin{multline*}
  \lexists[y_0]\dots\lexists[y_{k-1}]
  (!A_{g_0}(x_0,\dots,x_{l-1},y_0) \land
      \dots \land {}\\
  !A_{g_{k-1}}(x_0,\dots,x_{l-1},y_{k-1}) \land !A_f(y_0,\dots,y_{k-1},z))
\end{multline*}
\sourcecorrection{OLINC-027}{સ્થિર અંગ્રેજી
મૂળના આ પ્રદર્શિત સૂત્રમાં
પહેલો પરિમાણક $y_0\dots$ને
એક જ ચલ તરીકે લે છે અને
છેલ્લો સંયોજક અંદરના પરિમાણકની
હદ બહાર રહી જાય છે. અહીં
$y_0$થી~$y_{k-1}$ સુધીના
અસ્તિત્વ-પરિમાણકો તથા આખા
સંયોજનની હદ સ્પષ્ટ કરી છે.}
આ વિધેયને નિરૂપે છે:
\[
h(x_0, \dots, x_{l-1}) = f(g_0(x_0, \dots, x_{l-1}), \dots, g_{k-1}(x_0,
\dots, x_{l-1})).
\]
\end{prop}

\begin{proof}
અભ્યાસપ્રશ્ન.
\end{proof}

\begin{prob}
\olref[inc][req][cmp]{prop:rep1}
અને \olref[inc][req][cmp]{prop:rep2}ની
સાબિતીઓને માર્ગદર્શક બનાવી
\olref[inc][req][cmp]{prop:rep-composition}ની
સાબિતી વિગતે પૂર્ણ કરો.
\sourcecorrection{OLINC-028}{સ્થિર અંગ્રેજી
મૂળના આ અભ્યાસપ્રશ્નમાં
એક જ વિધાનનો સંદર્ભ બે વાર છે;
અગાઉની બે પૂરક સાબિતીઓમાંથી
પહેલી અને બીજી બંનેનો સંદર્ભ
અહીં આપ્યો છે.}
\end{prob}

\end{document}
